\section{Diagnósticos de ajuste: escolha dos moldes, gráficos QQ e
resíduos quantílicos}
\label{ap:diagnosticos}

\setcounter{figure}{0}
\renewcommand{\thefigure}{A.\arabic{figure}}
\setcounter{table}{0}
\renewcommand{\thetable}{A.\arabic{table}}

Este apêndice reúne todo o material estatístico de ajuste do documento
--- a escolha da família de distribuição de cada indicador, as notas de
qualidade de ajuste, os gráficos QQ e os resíduos quantílicos dos
modelos ---, para que o corpo do texto flua sem interrupções técnicas.
Ele é dirigido a quem quiser conferir a base metodológica; a leitura
substantiva não depende dele. Três noções bastam para acompanhar tudo.
\textbf{Nota AIC}: uma medida de qualidade de ajuste usada para comparar
moldes de distribuição entre si; quanto menor o AIC, melhor o molde
descreve os dados, e o valor isolado não tem interpretação --- só serve
para comparar. \textbf{Gráfico QQ}: compara os dados reais com um molde
matemático de distribuição; se os pontos seguem a reta, o molde descreve
bem os dados, e desvios sistemáticos da reta indicam onde o molde falha.
O número r que acompanha cada QQ resume essa aderência: quanto mais
próximo de 1, melhor. \textbf{Resíduo quantílico}: é a versão do QQ
para um modelo já ajustado com covariáveis; se o modelo descreve bem
os dados, esses resíduos se comportam como uma Normal padrão e os
pontos deitam na diagonal.

\subsection*{Escolha do molde de cada indicador}

Faturamento por saída e TMP são valores positivos com cauda à direita, e
\textbf{cinco moldes clássicos} de distribuição foram postos à prova
para cada um, comparados pela nota AIC e pelo gráfico QQ. O resultado é
consistente (Tabela~\ref{tab:familias}): o molde \textbf{log logístico
(também chamado de Fisk)} vence nas duas amostras, com a \textbf{LogNormal}
logo atrás, e a Gama, molde usado por tradição em análises de custo,
fica sempre entre as piores. No TMP o ajuste do log logístico é
praticamente perfeito (r de 0,9974 na Figura~\ref{fig:qqtmpsemlp}),
porque a consolidação do painel já removeu os hospitais de longa
permanência que, na versão anterior, criavam um teto de 30,5 dias que
nenhum molde contínuo reproduzia. Tanto no faturamento (r de 0,9700,
Figura~\ref{fig:qqcusto}) quanto no TMP (Figura~\ref{fig:qqtmp}) o molde
log logístico acompanha os dados por quase toda a extensão.

\begin{table}[H]
\centering
\caption{Comparação dos cinco moldes de distribuição para faturamento e TMP
pela nota AIC e pelo r do gráfico QQ.}
\label{tab:familias}
\small
\begin{tabular}{llrrr}
\toprule
Amostra & Molde & $n$ & AIC & r do QQ \\
\midrule
Faturamento por saída & Log logística (Fisk) & 3.179 & 50.785 & 0,9700 \\
                & LogNormal            & 3.179 & 50.883 & 0,9079 \\
                & Gaussiana inversa    & 3.179 & 50.952 & 0,8943 \\
                & Gama                 & 3.179 & 51.867 & 0,8310 \\
                & Weibull              & 3.179 & 52.320 & 0,8343 \\
\midrule
TMP & Log logística (Fisk) & 3.178 & 11.878 & 0,9974 \\
                     & LogNormal            & 3.178 & 11.956 & 0,9810 \\
                     & Gaussiana inversa    & 3.178 & 11.989 & 0,9786 \\
                     & Gama                 & 3.178 & 12.150 & 0,9631 \\
                     & Weibull              & 3.178 & 12.920 & 0,9398 \\
\bottomrule
\end{tabular}
\end{table}

\comofeito{cada bloco de linhas é uma amostra (faturamento por saída e
TMP) e cada linha, um molde de distribuição ajustado a ela. A coluna
\textbf{$n$} é o número de observações usadas; \textbf{AIC} é a nota de
ajuste obtida do ajuste de máxima verossimilhança de cada molde (quanto
menor, melhor); e \textbf{r do QQ} é a correlação entre os pontos do
gráfico QQ e a reta ideal (quanto mais perto de 1, melhor). Em cada
amostra os moldes aparecem ordenados do melhor para o pior AIC.}

Para a fração de alta complexidade, a Figura~\ref{fig:betapct} mostra o
ajuste Beta do miolo (os hospitais que realizam alguma alta
complexidade); os dois acúmulos nas pontas --- zero e um --- são
tratados separadamente pelo modelo ZOIB citado na
Seção~\ref{sec:distribuicoes}.

\comoler{a grande maioria dos hospitais que fazem alta complexidade faz
pouca (a concentração à esquerda; metade fica abaixo de 3,7\%), e uma
minoria é altamente especializada (as barras à direita, perto de
100\%). A curva escura é o molde Beta, que acompanha essa queda
contínua.}

\begin{figure}[H]
\centering
\includegraphics[width=.62\textwidth]{fig_ae_02_beta_pct_alta_complex.png}
\caption{Fração de alta complexidade entre os hospitais que realizam
alguma alta complexidade: histograma e molde Beta ajustado. O eixo
vertical está em escala logarítmica para a cauda inteira ficar
visível.}
\label{fig:betapct}
\end{figure}

\comofeito{o eixo horizontal é a fração de alta complexidade, restrita
aos hospitais que realizam alguma (entre 0 e 1), e o vertical acompanha
a quantidade de observações por faixa, em escala logarítmica (para a
cauda inteira ficar visível).
Sobre o histograma foi desenhada a curva do molde Beta ajustado a esses
dados.}

\subsection*{Gráficos QQ da fase exploratória}

As Figuras~\ref{fig:qqcusto} a~\ref{fig:qqtmpsemlp} avaliam os moldes
marginais de faturamento e TMP da Seção~\ref{sec:distribuicoes}.

\comoler{quanto mais os pontos abraçam a reta, melhor o molde. No
painel (b), o molde log logístico segue os dados por quase toda a
extensão, descolando apenas nos pouquíssimos hospitais mais caros; por
isso ele recebe a melhor nota da Tabela~\ref{tab:familias}.}

\begin{figure}[H]
\centering
\subfig{fig_ae_03_qq_lognormal_custo_saida.png}{Molde LogNormal}
\subfig{fig_ae_03_qq_fisk_custo_saida.png}{Molde log logístico (Fisk)}
\caption{Gráficos QQ do faturamento por saída: os pontos são os dados e a
reta é o comportamento esperado pelo molde.}
\label{fig:qqcusto}
\end{figure}

\comofeito{gráfico QQ. O eixo horizontal traz os quantis teóricos do
molde de distribuição e o vertical, os valores do faturamento por saída
ordenados do menor ao maior; a reta é o comportamento esperado se o
molde descrevesse os dados perfeitamente. Cada painel testa um molde
(LogNormal e log logístico/Fisk), e o r no título resume a aderência.}

\comoler{nos dois painéis os pontos acompanham a reta em quase toda a
extensão e descolam apenas nos pouquíssimos hospitais de permanência
mais longa. O painel corrente já não contém a subpopulação de longa
permanência que, na versão anterior desta análise, criava um teto de
30,5 dias e o característico achatamento no alto: a ETAPA F a removeu, e
o ajuste do TMP passou a ser quase perfeito.}

\begin{figure}[H]
\centering
\subfig{fig_ae_03_qq_lognormal_tmp.png}{Molde LogNormal}
\subfig{fig_ae_03_qq_fisk_tmp.png}{Molde log logístico (Fisk)}
\caption{Gráficos QQ do TMP no painel completo.}
\label{fig:qqtmp}
\end{figure}

\comofeito{mesma construção da figura anterior (quantis teóricos no eixo
horizontal, valores ordenados no vertical, reta de referência), agora
para o TMP no painel completo, comparando os moldes LogNormal e log
logístico.}

\comoler{este é o gráfico QQ do TMP no painel atual: os pontos seguem a
reta e o molde descreve os dados quase perfeitamente. Não há mais o
achatamento que a versão anterior do painel exibia, porque os hospitais
de longa permanência que o causavam saíram na consolidação.}

\begin{figure}[H]
\centering
\includegraphics[width=.62\textwidth]{fig_ae_03_qq_fisk_tmp_sem_longa_perm.png}
\caption{Gráfico QQ do TMP no painel de 289 hospitais (log logística):
os pontos ficam praticamente sobre a reta (r de 0,9974).}
\label{fig:qqtmpsemlp}
\end{figure}

\comofeito{gráfico QQ do TMP no painel de 289 hospitais, molde log
logístico: quantis teóricos no eixo horizontal, valores ordenados do TMP
no vertical, reta de referência. Quanto mais os pontos ficam sobre a
reta, melhor o ajuste (aqui r de 0,9974).}

\subsection*{Os moldes melhoram quando se compara o comparável}

Uma pergunta natural: se os QQ do painel completo não são perfeitos, o
molde está errado? As Figuras~\ref{fig:qqcondA} a~\ref{fig:qqcondC}
respondem que não: o desajuste vem de misturar 289 hospitais muito
diferentes num único gráfico. Quando se compara cada hospital consigo mesmo (controlando o
efeito do hospital e do ano), os desvios praticamente desaparecem: no
faturamento, a assimetria restante é quase nula; na mortalidade, a forte
assimetria da distribuição bruta na escala \textbf{logito} (a escala
interna em que o modelo trabalha proporções), de 1,97 negativa, cai
para 0,27 negativa depois do controle; na escala original da
Tabela~\ref{tab:univariada} a assimetria é outra (0,80 positiva), e os
dois números não se contradizem. Em resumo, \textbf{os moldes escolhidos descrevem bem os
dados quando se compara o comparável}.

\comoler{compare o gráfico de cima com o de baixo: em cima, todos os
hospitais juntos, os pontos serpenteiam em volta da reta; embaixo,
comparando cada hospital consigo mesmo, os pontos praticamente deitam
na reta.}

\begin{figure}[H]
\centering
\subfig{fig_ae_11_qq_custo_marginal.png}{Faturamento: todos os hospitais
juntos}
\subfig{fig_ae_11_qq_custo_residuo_ef.png}{Faturamento: cada hospital
comparado consigo mesmo}
\caption{Faturamento por saída: gráfico QQ antes e depois de controlar as
diferenças fixas entre hospitais e entre anos.}
\label{fig:qqcondA}
\end{figure}

\comofeito{dois gráficos QQ do faturamento (quantis teóricos no eixo
horizontal, valores ordenados no vertical). No painel de cima, todos os
289 hospitais entram juntos; no de baixo, controla-se o efeito fixo de
hospital e de ano, isto é, cada hospital é comparado consigo mesmo ao
longo do tempo. Comparar os dois isola se o desajuste vem do molde ou da
mistura de hospitais diferentes.}

\comoler{no painel anterior este par contrastava o desenho em degrau
(com os hospitais de longa permanência) com a reta quase perfeita (sem
eles). No painel corrente a ETAPA F já removeu essa subpopulação, então
as duas versões coincidem e ambas mostram a reta --- a confirmação de
que o desvio era a subpopulação, não o molde.}

\begin{figure}[H]
\centering
\subfig{fig_ae_11_qq_tmp_marginal.png}{TMP: todos os hospitais juntos}
\subfig{fig_ae_11_qq_tmp_marginal_sem_longa_perm.png}{TMP: sem a antiga
subpopulação de longa permanência}
\caption{TMP: gráfico QQ do painel completo e da versão sem a antiga
subpopulação de longa permanência --- no painel corrente as duas
coincidem, porque a ETAPA F já removeu essa subpopulação.}
\label{fig:qqcondB}
\end{figure}

\comofeito{mesma construção da figura anterior, para o TMP: o painel de
cima traz o painel completo e o de baixo, a versão sem a antiga
subpopulação de longa permanência (no painel atual as duas coincidem,
pois a ETAPA F já removeu essa subpopulação).}

\comoler{mesmo padrão do faturamento: a curva em S do gráfico de cima, com
todos os hospitais misturados, vira uma reta bem comportada quando cada
hospital é comparado consigo mesmo. A heterogeneidade entre hospitais,
e não o molde, era a fonte do desajuste.}

\begin{figure}[H]
\centering
\subfig{fig_ae_11_qq_mort_logito_marginal.png}{Mortalidade: todos os
hospitais juntos}
\subfig{fig_ae_11_qq_mort_logito_residuo_ef.png}{Mortalidade: cada
hospital comparado consigo mesmo}
\caption{Mortalidade: gráfico QQ antes e depois de controlar as
diferenças fixas entre hospitais e entre anos.}
\label{fig:qqcondC}
\end{figure}

\comofeito{mesma construção (QQ com todos os hospitais juntos em cima e
cada hospital comparado consigo mesmo embaixo), agora para a mortalidade
na escala logito, a escala interna em que o modelo trabalha proporções.}

\subsection*{Resíduos quantílicos dos modelos estimados}

A fase de estimação prometeu que a escolha das famílias de distribuição
seria validada nos modelos ajustados, e é o que as
Figuras~\ref{fig:qqprimeira} e~\ref{fig:qqsegunda} fazem, com a
Figura~\ref{fig:qqgama} documentando a reestimação da UTI comentada na
Seção~\ref{sec:inferenciafinal}. O resíduo quantílico é um teste de
qualidade: se o modelo descreve bem os dados, esses resíduos se alinham
sobre a diagonal do gráfico. O alinhamento é muito bom na mortalidade
(correlação de 0,976), na fração de alta complexidade (0,991) e na
produção (0,967; o piso achatado na ponta esquerda do QQ, com
observações empilhadas perto de menos 5, corresponde a hospitais
produzindo muito abaixo do que porte e perfil preveriam, desvio
confinado à cauda extrema que não compromete o contraste entre
categorias), e bom no faturamento (0,970) e no tempo de permanência
(0,952), cujas caudas pesadas já estavam diagnosticadas e são absorvidas
pelas margens de erro agrupadas. A exceção honesta é a
\textbf{intensidade de ocupação da UTI} (0,814): a cauda de UTIs com uso
quase nulo escapa do modelo em logaritmo, e esse resultado específico
deve ser lido com reserva; fica registrado para a próxima iteração o uso
de uma família com massa perto de zero para esse indicador --- reestimada
com a família Gama na Figura~\ref{fig:qqgama}.

\comoler{quanto mais os pontos azuis acompanham a linha tracejada,
melhor o modelo descreve os dados; o valor de r no canto é essa
aderência resumida em um número (1 seria perfeito). Compare o
alinhamento quase exato da fração de alta complexidade com o
descolamento da cauda esquerda na UTI, que motivou a reestimação da
figura seguinte.}

\begin{figure}[H]
\centering
\subfig{fig_est_04_qq_mort_zib.png}{Mortalidade geral, modelo Beta com
inflação em zero}
\subfig{fig_est_04_qq_pct_alta_zoib.png}{Fração de alta complexidade,
modelo ZOIB}
\caption{Resíduos quantílicos dos modelos de taxa.}
\label{fig:qqprimeira}
\end{figure}

\comofeito{gráficos de resíduos quantílicos de dois modelos já
estimados (mortalidade, Beta com inflação em zero; alta complexidade,
ZOIB). O eixo horizontal traz os quantis teóricos da Normal e o
vertical, os resíduos quantílicos ordenados; se o modelo descreve bem os
dados, os pontos deitam sobre a diagonal. O r no canto resume a
aderência.}

\begin{figure}[H]
\centering
\subfig{fig_est_04_qq_producao_bn.png}{Produção, Binomial Negativa}
\subfig{fig_est_04_qq_ocup_uti_pos.png}{Ocupação de UTI entre os
positivos, o único ajuste insatisfatório}
\caption{Resíduos quantílicos da produção e da parte positiva da UTI.}
\label{fig:qqsegunda}
\end{figure}

\comofeito{mesma construção da figura anterior (quantis teóricos da
Normal no eixo horizontal, resíduos quantílicos ordenados no vertical,
diagonal de referência), para os modelos de produção (Binomial Negativa)
e da parte positiva da ocupação de UTI.}

\comoler{estes são os quatro modelos restantes, cujos valores de
aderência o texto acima cita: faturamento e tempo de permanência na
especificação de efeitos fixos, ocupação de internação e a versão
ajustada da mortalidade. Em todos, os pontos acompanham a diagonal na
maior parte do trajeto, descolando apenas nas caudas.}

\begin{figure}[H]
\centering
\subfig{fig_est_04_qq_custo_fe.png}{Faturamento real por saída, efeitos
fixos}
\subfig{fig_est_04_qq_tmp_fe.png}{Tempo de permanência, efeitos fixos}
\caption{Resíduos quantílicos dos dois modelos em logaritmo com efeitos
fixos de hospital.}
\label{fig:qqterceira}
\end{figure}

\comofeito{mesma construção dos resíduos quantílicos anteriores (quantis
teóricos da Normal no eixo horizontal, resíduos ordenados no vertical,
diagonal de referência), para os modelos de faturamento real por saída e
de tempo de permanência, ambos estimados com efeitos fixos de hospital.}

\begin{figure}[H]
\centering
\subfig{fig_est_04_qq_ocup_internacao.png}{Ocupação de internação}
\subfig{fig_est_04_qq_mort_ajustada_zib.png}{Mortalidade ajustada,
modelo Beta com inflação em zero}
\caption{Resíduos quantílicos da ocupação de internação e da versão
ajustada da mortalidade.}
\label{fig:qqquarta}
\end{figure}

\comofeito{mesma construção, para o modelo de ocupação de internação e
para a versão da mortalidade que exclui os óbitos fetais, maternos e
neonatais --- esta última é a que sustenta a verificação de robustez da
Seção~\ref{sec:robustez}.}

\comoler{compare com o painel (b) da figura anterior: os pontos que
antes despencavam da linha tracejada na ponta esquerda agora a
acompanham em quase todo o trajeto.}

\begin{figure}[H]
\centering
\includegraphics[width=.5\textwidth]{fig_est_08_qq_uti_gama.png}
\caption{Resíduos quantílicos da intensidade de UTI reestimada com a
família Gama: a aderência sobe de 0,814 para 0,944.}
\label{fig:qqgama}
\end{figure}

\comofeito{mesma construção do resíduo quantílico (quantis teóricos da
Normal no eixo horizontal, resíduos ordenados no vertical), para a
intensidade de ocupação de UTI reestimada com a família Gama; a aderência
sobe de 0,814 para 0,944 em relação ao modelo em logaritmo.}

